\clearpage
\section*{Duplicate eLMS records made stalls look more common after 2019; most refiled stalls pass}
\textit{September 28, 2026. Agent-drafted exploratory work on branch \texttt{explore-alderman-measures}, after Jacob
asked how the stall rate was built and which choices behind it he had not seen.}

\textit{Duplicate records.} When the City Clerk moved to its new legislative system in 2023, each zoning map
amendment still pending was given a second record, numbered O2023-000xxxx and carrying the original introduction
date, on which its later actions were recorded; the original record stayed at ``Council Introduction.'' Counted
record by record, those originals looked like applications that lapsed without a vote. Jacob chose to fix this in
the producer: records introduced the same day with the same title or the same numeric application number are now one
amendment, represented by the record on which it was resolved. That leaves 6,568 amendments, of which 148 had two
records (67 passed, 4 pending, 77 stalled). Among applications with an alderman decided before the current term
(May 15, 2023), 15.2 percent of those introduced in 2019--2023 appear stalled when counted by record and 9.4 percent
when counted by amendment, against 6.5 percent in 2015--2018 and 7.4 percent in 2010--2014. The apparent doubling of
stalls after 2019 was mostly duplicates; a smaller rise remains. Re-geocoding the corrected amendments left all
5,800 addresses, wards and aldermen unchanged.

\textit{Refilings.} Many stalls are filed again. A later amendment is taken to refile an earlier one if the same
kind of filer submits it within four years with the same application number or an overlapping street segment on the
same side of the street; every chosen pair was read by hand, and 18 of 112 reviewed pairs were rejected as different
projects (mostly a single site against a corridor rezoning around it). 94 amendments were refiled, a median of 553
days later, and 85 of the refilings passed. Of the 318 stalls among the 4,102 decided applications, 78 were refiled
and passed, 5 were refiled without passing and 235 were never refiled. The zoning measures now separate terminal
stalls from delays by refiling, and, since Jacob asked for a minimum number of applications, report shrunk scores
only for aldermen with at least 20 decided applications (57 aldermen). Shrunk stall rates spread by 3.3 percentage
points across aldermen, but the earlier and later halves of an alderman's applications agree only at 0.23.

\textit{Project size.} Walter Burnett, Jr.\ has the longest adjusted days to passage; Jacob reads this as a sign
that the measures still reflect the projects a ward receives, since Burnett is broadly favorable to housing. To
control for project size, the applications' forms are now read for lot size, dwelling units, parking, commercial
space and height (\path{tasks/extract_zoning_application_forms}); against 60 forms read by hand, the parser reads
the lot size correctly for 51 and dwelling units for 39 of the 45 that state them, and stops if it contradicts an
unambiguous hand read. On 40 other forms, held out from its writing, it reads the lot size correctly for 34 of 37,
dwelling units for 17 of 28 and parking for 22 of 34; 8 of its 20 differences are fixable errors. The size controls
have not yet been added. Exploration outputs computed before the duplicate correction (boundary counts, placebo
boundaries, stall trends) are stale until rerun.

\textit{Reproduction.} Run \texttt{make} in \path{tasks/clean_zoning_map_amendments/code/},
\path{tasks/link_zoning_refilings/code/} and \path{tasks/estimate_alderman_zoning_measures/code/}; the record-count
comparison adds back one stalled record for each merged amendment that passed or stalled.
